\makeatletter\if@firstcolumn\null\newpage\fi\makeatother\newpage % Method starts on a new page
\section{Method}
\label{sec:methods}

\subsection{Overview of the pipeline}
\Cref{fig:pipeline} shows how the parts of the project fit together.
Both solvers start from an initial spin and write their results to CSV files, and the analysis and animation scripts work from these files.
\begin{figure}[htb]
    \centering
    \begin{tikzpicture}[node distance=0.5cm and 0.6cm,
        box/.style={draw,rounded corners=2pt,align=center,font=\scriptsize\sffamily,minimum height=0.9cm,text width=2.1cm},
        arr/.style={-{Stealth[length=1.8mm]},thick}]
        \node[box] (ic) {Initial spin\\ $\vb*{\omega}_0$};
        \node[box,right=of ic] (eng) {Spring-mass engine\\ Python / Rust, Verlet};
        \node[box,below=of ic] (rb) {Rigid-body solver\\ Euler eqs.\ + quaternions, RK4};
        \node[box,below=of eng] (csv) {CSV files};
        \node[box,below=of csv] (ana) {Analysis\\ NumPy / SciPy};
        \node[box,left=of ana] (out) {Plots \&\\ Manim videos};
        \draw[arr] (ic) -- (eng); \draw[arr] (eng) -- (csv); \draw[arr] (csv) -- (ana);
        \draw[arr] (ic) -- (rb); \draw[arr] (rb) -- (csv); \draw[arr] (ana) -- (out);
    \end{tikzpicture}
    \caption{Overview of the simulation and analysis pipeline. The spring engines start from $\vb{v}_i=\vb*{\omega}_0\times\vb{r}_i$, the rigid-body solver from $\vb*{\Omega}_0=\vb*{\omega}_0$ and $q_0=(1,\vb{0})$.}
    \label{fig:pipeline}
\end{figure}

\subsection{The rigid body as point masses and springs}
\label{sec:spring-model}
In our first approach no rigid body equations are coded at all.
The body is built from a few point masses, and every pair of masses is joined by a stiff spring whose rest length is the initial distance between them (\code{rigid=True} in \code{System}).
With $N$ masses there are $N(N-1)/2$ springs, so the six-mass octahedron has 15.
Each spring pulls along the line joining its two masses, with a force proportional to its stretch plus an optional damping term proportional to the rate of stretching (\cref{lst:spring}).
These forces are equal, opposite and central, so the total momentum and angular momentum are conserved, and damping only removes energy.
Rigid body motion, Euler's equations and the flip are not put in by hand; they have to emerge from Newton's second law for the individual masses.

\medskip\noindent\begin{minipage}{\linewidth}
\begin{minted}[bgcolor=LightGray,breaklines,fontsize=\scriptsize]{python}
def _spring_force(self, x, v, t):
    F = np.zeros_like(x)
    for (i, j, k, c, rest) in self.springs:
        d = x[j] - x[i]  # points from i to j
        r = np.linalg.norm(d)
        s = k*(r - rest) + c*np.dot(v[j] - v[i], d)/r
        f = s*d/r
        F[i] += f
        F[j] -= f
    return F
\end{minted}
\captionof{listing}{Spring forces on all masses (\code{engine.py}, condensed).}
\label{lst:spring}
\end{minipage}\par\medskip

The springs have to be stiff so that the body barely deforms.
With $k=\SI{e5}{N/m}$ and \SI{1}{kg} masses they vibrate at $\omega_s\sim\sqrt{k/m}\approx\SI{300}{rad/s}$, while the body itself turns at about \SI{1}{rad/s}.
The time step has to resolve these fast vibrations, which makes the model expensive.
For this reason the engine, first written in Python, was ported with the same structure to Rust for the long runs with small time steps.

\subsubsection{Velocity Verlet}
\label{sec:verlet}
Newton's equations for the masses are integrated with the velocity Verlet scheme in its kick--drift--kick form,
\begin{equation}
\begin{aligned}
    \vb{v}_i^{n+1/2} &= \vb{v}_i^{n} + \tfrac{\Delta t}{2m_i}\vb{F}_i(\vb{x}^{n})\\
    \vb{x}_i^{n+1}   &= \vb{x}_i^{n} + \Delta t\,\vb{v}_i^{n+1/2}\\
    \vb{v}_i^{n+1}   &= \vb{v}_i^{n+1/2} + \tfrac{\Delta t}{2m_i}\vb{F}_i(\vb{x}^{n+1})
\end{aligned}
\label{eq:verlet}
\end{equation}
Each mass gets half a velocity kick from the current forces, moves for a full step, and gets the second half kick from the forces at its new position.
The forces from the end of one step are reused at the start of the next, so only one force evaluation is needed per step (\cref{lst:verlet}).
Verlet is second order, time-reversible and symplectic, so the energy error stays bounded over long runs instead of drifting \cite{verlet1967,hairer2006}.
This is the very reason we chose to use this over RK4, which is more accurate over one step.

\medskip\noindent\begin{minipage}{\linewidth}
\begin{minted}[bgcolor=LightGray,breaklines,fontsize=\scriptsize]{python}
def step(self, dt):
    self.v += 0.5*dt*self.force_cached/self.m  # half kick
    self.x += dt*self.v                         # drift
    self.t += dt
    self.force_cached = self.total_force()
    self.v += 0.5*dt*self.force_cached/self.m  # half kick
\end{minted}
\captionof{listing}{One velocity Verlet step (\code{engine.py}; \code{RigidBody::step} in Rust is the same).}
\label{lst:verlet}
\end{minipage}\par\medskip

\subsection{Direct rigid body solver with quaternions}
\label{sec:quaternions}
While the spring model is general, it is indirect and expensive: the body is only approximately rigid, and the stiff springs force a tiny time step.
A truly rigid body needs much less information.
Its state is just its angular velocity and its orientation, and nothing in it vibrates, so the time step only has to resolve the rotation itself.
As an independent check we therefore also solve the rigid body equations directly (\code{quaternion.py}).
The angular velocity $\vb*{\Omega}$ in the body frame follows Euler's equations (\cref{eq:euler}).
The remaining question is how to store the orientation.

\paragraph{Choosing a representation}
Three Euler angles are the most familiar choice and use the minimum number of variables, but their equations of motion contain $1/\sin\theta$ and blow up when two rotation axes line up (this is called gimbal lock), and every step needs trigonometric functions.
A $3\times3$ rotation matrix has to store nine numbers for three degrees of freedom; numerical errors slowly make it non-orthogonal, so it stops being a pure rotation and has to be re-orthogonalised.
We found that there is a very beautiful way of describing these rotations using quaternions \cite{3b1b_quat}. These are complex numbers but in 4D.
A rotation by an angle $\phi$ about a unit axis $\vu{n}$ is written as
\begin{equation}
    q = \left(\cos\tfrac{\phi}{2},\ \vu{n}\sin\tfrac{\phi}{2}\right), \qquad \abs{q}=1,
    \label{eq:quat-def}
\end{equation}
a scalar and a vector part, and two quaternions are multiplied as
\begin{equation}
\begin{aligned}
    (a,\vb{a})\otimes(b,\vb{b}) = \big(&ab-\vb{a}\vdot\vb{b},\\
                                       &a\vb{b}+b\vb{a}+\vb{a}\times\vb{b}\big).
\end{aligned}
    \label{eq:quat-mult}
\end{equation}
Combining two rotations is a single product, with 16 multiplications instead of the 27 of a matrix product.

\paragraph{Equations of motion}
With this representation the orientation evolves as
\begin{equation}
    \dot q = \tfrac12\, q\otimes(0,\vb*{\Omega}),\qquad
    \vb{r}_{\text{lab}} = q\otimes(0,\vb{r}_{\text{body}})\otimes q^{*}.
    \label{eq:quat}
\end{equation}
The first equation is linear in $q$ and contains only products and sums, with no trigonometric functions and no singular points, which makes it both cheap and robust to integrate.
The only constraint is $\abs{q}=1$, which numerical errors break slowly and which is restored by simply dividing $q$ by its length.
The second relation rotates any point of the body into the lab frame, which we use for plots and animations.
In total the state is seven numbers, $\vb*{\Omega}$ and $q$, and \cref{eq:euler,eq:quat} are solved together (\cref{lst:derivs}).

\begin{sloppypar} % units and reference names would otherwise stick out of the column
\paragraph{The racket}
We use this solver for a real racket, the Babolat Pure Drive 2025, with the specifications listed by Tennis Warehouse \cite{tw_puredrive,twu_analyzer}: mass \SI{318}{g} strung, balance point \SI{32.99}{cm} from the butt, swingweight \SI{317}{kg.cm^2} and twistweight \SI{15.7}{kg.cm^2}.
For the racket we number the axes from the handle (1, smallest moment) to the face normal (3, largest), the reverse of \cref{sec:inertia}; Euler's equations (\cref{eq:euler}) hold in any right-handed principal frame, and \cref{eq:lambda} is unchanged when $I_1$ and $I_3$ are swapped.
The twistweight is the moment about the handle, $I_1=\SI{1.57e-3}{kg.m^2}$.
The swingweight is measured about an axis \SI{10}{cm} from the butt, so the parallel-axis theorem gives $I_2=\SI{1.49e-2}{kg.m^2}$ about the centre of mass, and because the racket is nearly flat the perpendicular-axis theorem gives $I_3=I_1+I_2=\SI{1.65e-2}{kg.m^2}$.
The racket is thrown with $\vb*{\Omega}_0=(-1.96,\,9.26,\,0)\,\si{rad/s}$, chosen so that after the \SI{0.7}{s} flight it has made exactly half a turn about its handle.
\end{sloppypar}

\medskip\noindent\begin{minipage}{\linewidth}
\begin{minted}[bgcolor=LightGray,breaklines,fontsize=\scriptsize]{python}
def euler_equations(w):
    w1, w2, w3 = w
    return np.array([(I2 - I3)/I1*w2*w3,
                     (I3 - I1)/I2*w3*w1,
                     (I1 - I2)/I3*w1*w2])

def dq(q, w):
    return 0.5*quaternion_multiply(q, [0, *w])
\end{minted}
\captionof{listing}{Time derivatives of $\vb*{\Omega}$ and $q$ (\code{quaternion.py}, condensed).}
\label{lst:derivs}
\end{minipage}\par\medskip

\subsubsection{Fourth-order Runge--Kutta}
\label{sec:rk4}
This system has no stiff springs, so we integrate it with the fourth-order Runge--Kutta method (RK4),
\begin{equation}
    \vb{y}_{n+1} = \vb{y}_n + \frac{\Delta t}{6}\left(\vb{k}_1+2\vb{k}_2+2\vb{k}_3+\vb{k}_4\right),
    \label{eq:rk4}
\end{equation}
with $\vb{y}=(\vb*{\Omega},q)$ and $\Delta t=\SI{1}{ms}$.
$\vb*{\Omega}$ and $q$ are advanced together, so each of the four stages evaluates both derivatives (\cref{lst:rk4}).
RK4 does not keep $\abs{q}=1$ exactly, so $q$ is renormalised after every step.

\medskip\noindent\begin{minipage}{\linewidth}
\begin{minted}[bgcolor=LightGray,breaklines,fontsize=\scriptsize]{python}
def rk4(w, q, dt):
    k1, l1 = euler_equations(w), dq(q, w)
    w2 = w + 0.5*dt*k1
    k2, l2 = euler_equations(w2), dq(q + 0.5*dt*l1, w2)
    w3 = w + 0.5*dt*k2
    k3, l3 = euler_equations(w3), dq(q + 0.5*dt*l2, w3)
    w4 = w + dt*k3
    k4, l4 = euler_equations(w4), dq(q + dt*l3, w4)
    w = w + dt/6*(k1 + 2*k2 + 2*k3 + k4)
    q = q + dt/6*(l1 + 2*l2 + 2*l3 + l4)
    return w, q/np.linalg.norm(q)  # keep |q| = 1
\end{minted}
\captionof{listing}{One RK4 step for $\vb*{\Omega}$ and $q$ (\code{quaternion.py}, condensed).}
\label{lst:rk4}
\end{minipage}\par\medskip

\subsection{Observables from particle positions}
\label{sec:observables}
The spring engines only save the positions of the masses, 60 times per simulated second.
Everything else is computed from these positions afterwards (\code{src/sim/}); only the conservation checks of \cref{sec:conservation} use the engine's own velocities.
\begin{sloppypar} % long words would otherwise stick out of the column
\begin{itemize}
    \item \textbf{Body axes.} The masses sit in opposite pairs on the principal axes, so the unit vector from one mass of a pair to the other gives that body axis $\vb{e}_k$.
    \item \textbf{Velocities.} Computed with \code{np.gradient}, which takes second-order finite differences of the saved positions.
    \item \textbf{Angular velocity.} Each axis moves as $\dot{\vb{e}}_k=\vb*{\omega}\times\vb{e}_k$; summing $\vb{e}_k\times\dot{\vb{e}}_k$ over the three axes gives $2\vb*{\omega}$, so
    \begin{equation}
        \vb*{\omega} = \frac12\sum_{k=1}^{3}\vb{e}_k\times\dot{\vb{e}}_k,
        \qquad \Omega_k = \vb*{\omega}\vdot\vb{e}_k.
        \label{eq:omega-from-axes}
    \end{equation}
    \item \textbf{Angular momentum.} $\vb{L}=\sum_i m_i\,\vb{r}_i\times\vb{v}_i$ about the centre of mass, projected onto the body axes to draw the polhode.
    \item \textbf{Flips.} A flip is counted each time the component of $\vb{L}$ along the intermediate axis changes sign.
\end{itemize}
\end{sloppypar}
For the quaternion solver $\vb*{\Omega}$ and $q$ are the state itself, so nothing has to be reconstructed.

\subsection{Implementation}
\label{sec:implementation}
\begin{sloppypar} % long code names would otherwise stick out of the column
The spring engine was first written in Python (\code{engine.py}) and then rewritten in Rust with the same design.
Both are small general-purpose engines: they know nothing about rigid bodies or Euler's equations, only about point masses and the forces between them.

\paragraph{Python engine}
A \code{System} object stores the positions, velocities and masses of all particles as NumPy arrays.
Masses are added with \code{add_mass}, springs with \code{add_spring}, and any other force (for example gravity or an external torque) with \code{add_force}, as a function that takes the current positions, velocities and time and returns the force on every mass.
With \code{rigid=True} the system joins every pair of masses with a spring before the first step.
At each step the forces from all these sources are summed and passed to the Verlet integrator (\cref{lst:verlet}), so adding new physics never requires touching the integrator.
The loop over springs runs in interpreted Python, which makes runs with $\Delta t=\SI{e-5}{s}$ ($10^7$ steps) impractically slow.
\end{sloppypar}

\paragraph{Rust engine}
The Rust version is a library split into small modules (\cref{lst:tree}).
\code{Vector3} is a plain 3D vector with overloaded arithmetic, dot and cross products, so the physics code reads like the equations.
A \code{RigidBody} owns its masses, springs and external forces, and its \code{step} is the same Verlet step as in Python.
An \code{Engine} holds one or more bodies and the clock, and each experiment is a short program in \code{examples/} that builds a body and runs it (\cref{lst:usage}).
Being compiled, the Rust engine is about 500 times faster than the Python one on the same run (\cref{tab:performance}): a $10^7$-step run takes under a second instead of seven minutes.

\medskip\noindent\begin{minipage}{\linewidth}
\begin{minted}[bgcolor=LightGray,breaklines,fontsize=\scriptsize]{text}
src/
  math/vector3.rs     Vector3 and its operators
  body/mass.rs        Mass {m, x, v}
  body/rigid_body.rs  springs, forces, Verlet step
  physics/spring.rs   damped spring force
  physics/force.rs    external force closures
  engine.rs           bodies, clock, CSV export
examples/             one program per run
\end{minted}
\captionof{listing}{Structure of the Rust engine (\code{physics-engine-rs}).}
\label{lst:tree}
\end{minipage}\par\medskip

\medskip\noindent\begin{minipage}{\linewidth}
\begin{minted}[bgcolor=LightGray,breaklines,fontsize=\scriptsize]{rust}
let mut b = RigidBody::new(true, 1e5, 0.0); // rigid, k, c
for r in positions {
    b.add_mass(m, r, w.cross(&r)); // v = w x r
}
let mut u = Engine::new();
u.add_body(b);
u.export_csv("six_masses_rs.csv", 1e-5, 60.0, 100.0)?;
\end{minted}
\captionof{listing}{Setting up and running a simulation with the Rust engine (condensed from \code{examples/}).}
\label{lst:usage}
\end{minipage}\par\medskip

\paragraph{Output and analysis}
Both engines write the same CSV format: a short header with $\Delta t$, the output rate, the masses and the list of springs, followed by one row per frame with the time and all positions.
The animation scripts in \code{src/sim/} therefore work with either engine.
They compute the observables of \cref{sec:observables} with NumPy \cite{numpy} and produce the Manim animations \cite{manim}.
We used AI tools to write the code that creates the videos and the code that exports the CSV files in this format after the runs.
Every other piece of code in the Python and Rust engines and the quaternion file is handwritten. We could not have made the beautiful videos demonstrating this effect without AI and consider this as fair use.


\begin{table}[htb]
    \centering
    \caption{Wall-clock time to integrate the octahedron for $T=\SI{100}{s}$ with $k=\SI{e5}{N/m}$, same initial state. After $10^5$ steps both engines end with bit-for-bit identical positions and velocities.}
    \label{tab:performance}
    \small
    \begin{tabular}{@{}lccc@{}}
        \toprule
        $\Delta t$ (s) & Steps & Python (s) & Rust (s) \\
        \midrule
        $10^{-3}$ & $10^5$ & 4.3 & 0.01 \\
        $10^{-5}$ & $10^7$ & 429 & 0.73 \\
        \bottomrule
    \end{tabular}
\end{table}
This shows how terribly slow Python is for computation whenever loops are involved.
